\uuid{P5j6}
\exo7id{3616}
\titre{exo7 3616}
\auteur{quercia}
\organisation{exo7}
\datecreate{2010-03-10}
\isIndication{false}
\isCorrection{false}
\chapitre{Réduction d'endomorphisme, polynôme annulateur}
\sousChapitre{Sous-espace stable}
\cpge{mpsi-mp2i.dimension-finie.c, pcsi.dimension-finie.c, ptsi.dimension-finie.c, mpsi-mp2i.applications-lineaires.b, pcsi.applications-lineaires.b, ptsi.applications-lineaires.b, mp-mpi.reduction-endomorphismes.b, pc.reduction.a, psi.reduction-endomorphismes.a, pt.reduction-endomorphismes.a, mp-mpi.reduction-endomorphismes.d, pc.reduction.c, psi.reduction-endomorphismes.c, pt.reduction-endomorphismes.c, mp-mpi.reduction-endomorphismes.i, pc.reduction.d, psi.reduction-endomorphismes.d}
\module{Algèbre}
\niveau{L2}
\difficulte{}

\contenu{
\question{
Un endomorphisme $f$ est dit semi-simple si tout sous-espace stable par~$f$
admet un supplémentaire stable par~$f$.
Montrer qu'un endomorphisme d'un $\C$-ev de dimension finie est semi-simple
si et seulement s'il est diagonalisable.
}
}
